\chapter{Derived Euler fibers and higher operations}
\label{df:chap:fibers}

The Euler quotient of Chapter~\ref{ed:chap:euler} is the cohomology
of a projective bundle. Its parameter acts by the first Chern class of
the tautological line. Taking the derived fiber at parameter zero
recovers the cochain algebra of the unit sphere bundle.
The resulting cohomology need not determine its
multiplication up to homotopy. For the tangent bundle of a K3 surface,
a nonzero triple Massey product records information lost by passing
to the cohomology algebra.

We shall construct the comparison with sphere-bundle cochains and compute
all the transferred products. We then retain the parameter action on
the Euler quotient. Its Jordan blocks give a second family of higher
operations, now on a module over an exterior algebra. Finally, split
bundle embeddings provide pullback and Gysin maps between the Euler
models. These maps exhibit the compatibility a geometric correspondence
must possess before it acts on a derived Euler fiber.

Two gradings enter these constructions. In the geometric cochain model,
the parameter $z$ has cohomological degree two. In the module calculation,
it has resolution degree zero and may retain degree two as a separate
internal grading. The unit in Theorem~\ref{ed:thm:factor} of
Chapter~\ref{ed:chap:euler} is generally
inhomogeneous in cohomological degree. Its equality of completed Euler
and elliptic ideals therefore does not itself give a comparison of
graded cochain algebras.

\section{The projective-bundle cochain algebra}
\label{df:sec:projective}

Let $X$ be a compact smooth manifold and let $E\to X$ be a complex
vector bundle of rank $r\geq1$. Write
$\pi:P=\mathbf P(E)\to X$ for its bundle of lines and
$L=\mathcal O_P(-1)\subset\pi^*E$ for the tautological line.
Over $\mathbb R$ use the de Rham algebra $A_X=\Omega^*(X)$.
Over $\mathbb Q$ use $A_X=A_{\mathrm{PL}}(X)$, the commutative
algebra of rational polynomial forms. Choose closed representatives
$c_j(E)\in A_X^{2j}$ of the Chern classes, with $c_0(E)=1$.

The Euler relation is the homogeneous polynomial
\begin{equation}\label{df:eq:euler}
 e_z(E)=\sum_{j=0}^r(-z)^{r-j}c_j(E),\qquad |z|=2.
\end{equation}
To impose this relation at the cochain level, adjoin one odd generator:
\begin{equation}\label{df:eq:projective}
 \mathcal P_E=A_X[z]\otimes\Lambda(\eta_E),\qquad
 |\eta_E|=2r-1,\qquad D\eta_E=e_z(E),\quad Dz=0.
\end{equation}
Here $\Lambda(\eta_E)$ is the exterior algebra, and $D$ extends the
differential of $A_X$ by the graded Leibniz rule. All differentials
in this section raise cohomological degree by one.

An explicit comparison over $\mathbb R$ determines the geometric meaning
of the generator. Choose Hermitian connections on $E$ and $L$, normalized
by
\[
 c(\nabla)=\det\left(1+\frac{i}{2\pi}F_\nabla\right),
 \qquad \lambda=c_1(\nabla L).
\]
The tautological inclusion is a nowhere-zero section of
$W=\pi^*E\otimes L^*$. Let $\nabla_0$ preserve the trivial line
spanned by this section and be trivial on it. Let $\nabla_1$ be the
tensor-product connection. Their top Chern forms satisfy
\[
 c_r(\nabla_0)=0,\qquad c_r(\nabla_1)=e_\lambda(\pi^*E).
\]

Put $a=\nabla_1-\nabla_0$ and $\nabla_t=\nabla_0+ta$.
If $F_t$ is the curvature and $P_r$ is the polarized invariant
polynomial normalized by $P_r(F,\ldots,F)=c_r(\nabla)$, define
\begin{equation}\label{df:eq:transgression}
 \Theta_E=r\int_0^1P_r(a,F_t,\ldots,F_t)\,dt.
\end{equation}
The Bianchi identity gives
$\frac{d}{dt}c_r(\nabla_t)=r\,dP_r(a,F_t,\ldots,F_t)$.
Consequently $d\Theta_E=e_\lambda(\pi^*E)$, and there is a
cochain-algebra map
\begin{equation}\label{df:eq:comparison}
 \Phi_E:\mathcal P_E\longrightarrow\Omega^*(P),\qquad
 \alpha\longmapsto\pi^*\alpha,\quad z\longmapsto\lambda,
 \quad\eta_E\longmapsto\Theta_E.
\end{equation}
This is the Chern--Simons transgression for the tautological
section.\footnote{D.~L.~Johnson, \emph{Chern--Simons forms on associated bundles,
and boundary terms},
\href{https://arxiv.org/abs/math/0601182}{math/0601182},
Theorem~1.2 and \S2.2.}

\begin{theorem}\label{df:thm:projective}
The map $\Phi_E$ is a quasi-isomorphism. Over $\mathbb Q$,
$\mathcal P_E$ is a commutative rational cochain model of
$\mathbf P(E)$, with $z$ representing $c_1(L)$.
\end{theorem}

\begin{proof}
As a complex, $\mathcal P_E$ is the extension defined by
multiplication by $e_z(E)$ on $A_X[z]$, with the shift in
\eqref{df:eq:projective}. This multiplication is injective on
$H^*(X)[z]$. Indeed, its leading coefficient is $(-1)^r$, so
the highest nonzero coefficient of a polynomial cannot disappear.
Zero divisors in $H^*(X)$ do not affect this argument. The long exact
sequence of the extension gives
\[
 H^*(\mathcal P_E)=H^*(X)[z]/(e_z(E)).
\]

The classes $1,c_1(L),\ldots,c_1(L)^{r-1}$ restrict to a basis on
each projective fiber. The projective-bundle theorem therefore identifies
this quotient with $H^*(P)$ by base pullback and $z\mapsto c_1(L)$.
These are exactly the maps induced by $\Phi_E$.

For rational polynomial forms, choose a closed representative of
$c_1(L)$. The top Chern relation is exact on $P$, so choose a primitive
as the image of $\eta_E$. The same cohomology calculation proves that
the resulting rational map is a quasi-isomorphism. If a representative
$e$ changes to $e+D\sigma$, the generator change
$\eta'_E\mapsto\eta_E+\sigma$ compares the two models.
\end{proof}

We use the polynomial model for the derived construction below.
Completion in $z$ means degreewise completion with $|z|=2$.
Its cohomology quotient agrees with the polynomial quotient because
$c_1(L)$ is nilpotent on $H^*(P)$. A degree-zero ring
$\mathbb Q[[q]]$ may be adjoined by completed scalar extension.

\section{The unit sphere bundle}
\label{df:sec:sphere}

Choose a Hermitian metric on $E$. A unit vector determines its line,
so the unit circle bundle of $L$ is the unit sphere bundle $S(E)$.
Thus $S(E)$ is the homotopy fiber of the classifying map
$P\to BU(1)$ of $L$. We write $C^*(-;\mathbb Q)$ for a commutative
rational cochain model, and give its parameter map
$\mathbb Q[z]\to C^*(P;\mathbb Q)$ by this classifying map.

\begin{theorem}\label{df:thm:sphere}
There is a multiplicative derived comparison
\begin{equation}\label{df:eq:derived-fiber}
 C^*(\mathbf P(E);\mathbb Q)
       \otimes^{\mathbb L}_{\mathbb Q[z]}\mathbb Q
 \simeq C^*(S(E);\mathbb Q),\qquad |z|=2.
\end{equation}
A cochain model for this fiber is
\begin{equation}\label{df:eq:sphere}
 \mathcal S_E=A_X\otimes\Lambda(\eta'_E),\qquad
 |\eta'_E|=2r-1,\qquad D\eta'_E=c_r(E).
\end{equation}
\end{theorem}

\begin{proof}
Resolve the augmentation of $\mathbb Q[z]$ by adjoining
$\nu$ of degree one with $D\nu=z$. The polynomial
\[
 Q_E(z)=\frac{e_z(E)-c_r(E)}{z}
\]
is closed and has degree $2r-2$. In the resulting relative model set
$\eta'_E=\eta_E-\nu Q_E(z)$. Its differential is $c_r(E)$.
The remaining factor $\mathbb Q[z]\otimes\Lambda(\nu)$ contracts
onto its constants. This computes the derived tensor product without
inverting the parameter.

The same change of generator has a geometric interpretation. On the
circle bundle put $\alpha=\frac{i}{2\pi}A_L$, where $A_L$ is the
unitary connection form. With the complex orientation of the circle,
$\int_{S^1}\alpha=-1$ and $d\alpha$ is the pullback of $\lambda$.
If $\rho:S(E)\to P$ is the projection, the form
\[
 \rho^*\Theta_E-\alpha Q_E(\rho^*\lambda)
\]
has differential equal to the pullback of $c_r(E)$ from $X$.

On a sphere fiber the first term vanishes by dimension. The second
is $(-1)^{r+1}\alpha\lambda^{r-1}$, with pullbacks suppressed.
Since $\int_{\mathbf P^{r-1}}\lambda^{r-1}=(-1)^{r-1}$, its integral
on the complex-oriented sphere is $-1$. It represents a fiber generator
with the sign fixed by $D\eta'_E=c_r(E)$. The sphere-bundle filtration,
whose fiber cohomology has two rows, now proves the de Rham comparison.

Over $\mathbb Q$ the relative model of an oriented odd-sphere bundle
adjoins a generator with differential its Euler class. The complex
structure makes its orientation local system trivial, and its Euler
class is $c_r(E)$. This gives \eqref{df:eq:sphere}. Equivalently,
the cochain comparison for the homotopy pullback over the simply
connected space $BU(1)$ gives \eqref{df:eq:derived-fiber}.
\end{proof}

\section{The multiplication retained by the fiber}
\label{df:sec:minimal}

Let $X$ now be a connected smooth projective complex variety of
dimension $d>0$, with $H^{\mathrm{odd}}(X;\mathbb Q)=0$.
Put $A=H^*(X;\mathbb Q)$ and assume that the topological Euler
characteristic $\chi=\chi(X)$ is nonzero. Choose
$\omega\in A^{2d}$ with $\int_X\omega=1$, and let
$\tau:A\to\mathbb Q$ be integration, extended by zero outside
degree $2d$. The top Chern class of $TX$ is $\chi\omega$.

A cochain algebra is \emph{formal} if it is connected to its cohomology
algebra with zero differential by a chain of quasi-isomorphisms.
Compact K\"ahler formality, including its rational form, permits the
replacement of $A_X$ by $A$.\footnote{P.~Deligne, P.~Griffiths,
J.~Morgan and D.~Sullivan, \emph{Real homotopy theory of K\"ahler
manifolds}, Inventiones Mathematicae \textbf{29} (1975), 245--274.
See
\href{https://www.math.stonybrook.edu/~dennis/publications/PDF/DS-pub-0024.pdf}
{\S6, Main Theorem}. Rational formality is proved by
D.~Sullivan, \emph{Infinitesimal computations
in topology}, Publications Math\'ematiques de l'IH\'ES
\textbf{47} (1977), 269--331.
See
\href{https://www.math.stonybrook.edu/~dennis/publications/PDF/DS-pub-0039.pdf}
{Theorem~12.1}.}

An exact change in the representative of the Euler class is absorbed
by the generator change used above. The sphere bundle $S(TX)$ therefore
has the finite cochain model
\begin{equation}\label{df:eq:finite}
 B=A\otimes\Lambda(\eta),\qquad
 |\eta|=2d-1,\qquad d\eta=\chi\omega.
\end{equation}

There is only one contractible pair, namely $\eta\mapsto\chi\omega$.
Set $A'=\bigoplus_{0<j<2d}A^j$. Representatives for cohomology form
the subspace
\begin{equation}\label{df:eq:cohomology}
 H=\mathbb Q1\oplus A'\oplus\eta A^{>0}\subset B.
\end{equation}
Let $i:H\to B$ be inclusion and let $p:B\to H$ kill $\eta$ and
$\omega$. Define the degree-minus-one homotopy $h$ by
\begin{equation}\label{df:eq:contraction}
 h(\omega)=\eta/\chi,\qquad h=0
 \quad\text{on the other summands}.
\end{equation}
Then $pi=1$, $dh+hd=1-ip$, and $h^2=hi=ph=0$.

The contraction transfers multiplication to higher maps on $H$.
To fix their signs, let $sV$ lower the degree of a graded space $V$ by one.
Give the tensor coalgebra
$T^c(sV)=\bigoplus_{n\geq0}(sV)^{\otimes n}$ its deconcatenation
coproduct. An uncurved $A_\infty$ structure on $V$ is a degree-one
coderivation $b$ with $b(1)=0$ and $b^2=0$.
Its Taylor components are the suspended multiplications.

We use
\[
 b_1(sa)=-s(da),\qquad
 b_2(sa,sb)=(-1)^{|a|}s(ab),\qquad
 b_3(sa,sb,sc)=(-1)^{2|a|+|b|}s\,m_3(a,b,c).
\]
Extending a Taylor component across preceding suspended inputs carries
their Koszul sign. A minimal structure has $b_1=0$, and a strict unit
acts as the ordinary unit in arity two and gives zero in higher arity.

\begin{theorem}\label{df:thm:minimal}
The space $H$ has a strictly unital minimal $A_\infty$ structure
quasi-isomorphic to $B$. Its binary product is $p(i(-)i(-))$.
The only possibly nonzero higher multiplication is
\begin{equation}\label{df:eq:ternary}
 m_3(a,b,c)=\frac{\tau(bc)a-\tau(ab)c}{\chi}\,\eta,
 \qquad a,b,c\in A'.
\end{equation}
It is zero on all other homogeneous input patterns, and $m_n=0$ for
$n\geq4$. The suspended components of a quasi-isomorphism $I:H\to B$
are
\begin{equation}\label{df:eq:inclusion}
 I_1=si s^{-1},\qquad
 I_2(sa,sb)=s\bigl(\tau(ab)\eta/\chi\bigr)
 \quad(a,b\in A'),\qquad I_n=0\quad(n\geq3),
\end{equation}
with $I_2$ zero on the remaining input patterns.
\end{theorem}

\begin{proof}
The projected binary product removes exactly the top-class component
$\tau(ab)\omega$. Its primitive is $h(ab)=\tau(ab)\eta/\chi$.
For $a,b\in A'$ the binary bar equation is therefore
\[
 -s\,d\bigl(\tau(ab)\eta/\chi\bigr)+s(ab)
       =s\bigl(ab-\tau(ab)\omega\bigr).
\]
The two placements of this primitive in arity three give
$p\bigl(a\,h(bc)-h(ab)\,c\bigr)$, which is
\eqref{df:eq:ternary}. If an input contains $\eta$, a prospective
contribution either has zero homotopy or contains $\eta^2$.
Normalization of the contraction gives strict unitality.

In a transferred multiplication tree, every nonzero homotopy edge
has output proportional to $\eta$. A second homotopy above it vanishes,
and the product of two such outputs vanishes. Every tree with at least
four leaves has one of these forms. The trees for $I_n$, $n\geq3$,
vanish for the same reason.

Write $b_B,b_H$ for the two bar coderivations. The remaining terms in
$b_BI=Ib_H$ are products of two $I_2$ outputs or applications of
$I_2$ to an $m_3$ output. Both are zero.
The linear term of $I$ is injective, hence so is $I$ by tensor length.
The equality $Ib_H^2=b_B^2I=0$ gives all the $A_\infty$ identities.

Since $I_1$ induces the
identity on cohomology, $I$ is a quasi-isomorphism.\footnote{B.~Keller,
\emph{Introduction to $A$-infinity algebras and modules},
\href{https://arxiv.org/abs/math/9910179}{math/9910179},
\S\S3.3--3.6 and 4.2. The displayed contraction fixes the signs here.}
\end{proof}

For example, the binary product of $a,b\in A'$ is
$ab-\tau(ab)\omega$, whereas $a(\eta b)=\eta ab$.
The product of two classes containing $\eta$ is zero.
The ternary product measures the difference between using the two
possible primitives when a triple product passes through the omitted
top class.

\section{Nonformality and examples}
\label{df:sec:nonformal}

If even cocycles $a,b,c$ have $ab=dU$ and $bc=dV$, the cocycle
$aV-Uc$ defines a triple Massey product. Its indeterminacy consists
of the products obtained by changing $U$ and $V$ by closed classes.
A nonzero Massey product with zero indeterminacy obstructs formality.
The contraction above supplies both primitives and their signs.

\begin{corollary}\label{df:cor:nonformal}
Under the hypotheses of Section~\ref{df:sec:minimal}, if
$\dim H^2(X;\mathbb Q)\geq2$, the rational cochain algebra of
$S(TX)$ is not formal. This applies to every $X=Y^{[n]}$ with $n>0$
and $Y$ a projective K3 or Enriques surface.
\end{corollary}

\begin{proof}
Choose independent $a,c\in A^2$. Poincar\'e duality gives
$b\in A^{2d-2}$ with $\tau(ab)=1$ and $\tau(bc)=0$.
Then $ab=\omega=d(\eta/\chi)$ and $bc=0$, so
\begin{equation}\label{df:eq:massey}
 \langle[a],[b],[c]\rangle=-[c\eta]/\chi\ne0.
\end{equation}
The indeterminacy is
$[a]H^{2d-1}(B)+H^{2d-1}(B)[c]$, which is zero.
Indeed, the base has no odd cohomology and $\eta$ is not closed.
The nonzero class $c\eta$ occurs in \eqref{df:eq:cohomology}.

The Hilbert schemes in the assertion are smooth projective.
In the Hilbert-scheme cohomology formula, a surface class of degree
$j$ contributes a generator of degree $j+2m-2$ in length $m$.
For these surfaces $j$ is even, and the degree-two surface classes
persist in every positive length.

The Euler generating series is
\[
 \sum_{n\geq0}\chi(Y^{[n]})p^n
       =\prod_{m\geq1}(1-p^m)^{-\chi(Y)}.
\]
Its coefficients are positive since $\chi(Y)$ is $24$ or $12$.
Thus all hypotheses hold. The cohomology formula used here is the
geometric Heisenberg decomposition\cite{NakajimaHilbert1997}.
\end{proof}

For a surface $Y$ satisfying the same hypotheses, write
$u_i u_j=Q_{ij}\omega$ in a basis of
$H^2(Y;\mathbb Q)$. The cohomology of $S(TY)$ has dimensions
$1,b_2(Y),b_2(Y),1$ in degrees $0,2,5,7$. Its products include
\begin{equation}\label{df:eq:surface-products}
 u_i(u_j\eta)=Q_{ij}\omega\eta,\qquad
 m_3(u_i,u_j,u_k)
       =\frac{Q_{jk}u_i-Q_{ij}u_k}{\chi(Y)}\,\eta.
\end{equation}
In a hyperbolic plane take $u^2=v^2=0$, $uv=\omega$,
$a=u+v$, and $c=u-v$. Then $a^2=2\omega$ and $ac=0$, giving
\[
 \langle a,a,c\rangle=-\frac{2}{\chi(Y)}c\eta
 =\begin{cases}
   -c\eta/6,&Y\text{ an Enriques surface},\\
   -c\eta/12,&Y\text{ a K3 surface}.
  \end{cases}
\]
Both Massey products have zero indeterminacy.

\begin{example}\label{df:ex:grassmannian}
The condition $b_2(X)\geq2$ is sufficient but not necessary.
Let $X=\operatorname{Gr}(2,4)$, the Grassmannian of two-dimensional
subspaces of $\mathbb C^4$, and let $S$ be its tautological rank-two
bundle. Put $h=c_1(S^*)$ and $p=c_2(S^*)$. Its rational cohomology is
\[
 A=\mathbb Q[h,p]/(h^3-2hp,\ h^4-3h^2p+p^2),
 \qquad |h|=2,\quad |p|=4.
\]

The Schubert classes are $1,h,p,h^2-p,hp,p^2$.
The relations are the degree-six and degree-eight terms of
$c(S)^{-1}$, because the quotient bundle has rank two.
They give $p^2=h^2p=\omega$, with $\int_X\omega=1$, and
$h^4=2\omega$. The six Schubert cells give $\chi(X)=6$.

Although $b_2(X)=1$, take $a=b=p$ and $c=h^2-p\ne0$.
Then $ab=\omega$ and $bc=0$. In the sphere model
$|\eta|=7$ and $d\eta=6\omega$, so
\[
 \langle[p],[p],[h^2-p]\rangle
       =-[(h^2-p)\eta]/6\ne0.
\]
Its indeterminacy vanishes because $H^7(B)=0$.
Thus $S(T\operatorname{Gr}(2,4))$ is also nonformal.
\end{example}

The general sphere-bundle formality criterion imposes a condition on
the full cohomology algebra of the base. A single Betti number cannot
replace it.\footnote{J.~Zhou, \emph{Formality of sphere bundles},
Algebraic \& Geometric Topology \textbf{25} (2025), 2297--2315,
\href{https://arxiv.org/abs/2304.09594}{2304.09594}, Theorem~1.2.
For a nonzero top Euler class, formality of the total space forces
the formal base cohomology to be generated by one element.}

\section{The higher action of the parameter}
\label{df:sec:parameter}

The preceding products belong to the cochain algebra of a sphere bundle.
To retain the action of $z$ on the projective-bundle quotient, we instead
take a derived Hom over the parameter ring. Set $R_0=\mathbb Q[[z]]$,
with $z$ in resolution degree zero, and let $C$ be a finite-length
$R_0$-module. Multiplication by $z$ is a nilpotent operator $N$ on $C$.
The Koszul dual algebra is
$\mathcal K=\operatorname{Ext}_{R_0}^*(\mathbb Q,\mathbb Q)
=\Lambda(\xi)$, where $\xi$ has resolution degree one.

A dg $\mathcal K$-module representing
$\operatorname{RHom}_{R_0}(\mathbb Q,C)$ is
\begin{equation}\label{df:eq:parameter-module}
 M_C=C\otimes\Lambda(\xi),\qquad
 d(a)=\xi Na,\quad d(\xi a)=0.
\end{equation}
The action of $\xi$ is exterior multiplication. Using the free resolution
of $\mathbb Q$ with differential $-z$ gives the displayed sign after
applying Hom. If $z$ has internal degree two as well, give $\xi$ internal
degree $-2$. Then $d$ raises resolution degree by one and preserves
internal degree.

For the cyclic block $C_m=R_0/(z^m)$ choose a basis
$e_0,\ldots,e_{m-1}$ with $Ne_j=e_{j+1}$ and $Ne_{m-1}=0$.
The cohomology of $M_{C_m}$ has one generator $u$ in resolution degree
zero and one generator $v$ in degree one. A contraction is
\begin{equation}\label{df:eq:block-contraction}
 \begin{aligned}
 i(u)&=e_{m-1},& i(v)&=(-1)^{m-1}\xi e_0,\\
 h(\xi e_j)&=e_{j-1}\quad(j>0),& h&=0\quad\text{otherwise}.
 \end{aligned}
\end{equation}
The projection is $p(e_{m-1})=u$, $p(\xi e_0)=(-1)^{m-1}v$,
and zero on the other basis vectors. It satisfies $dh+hd=1-ip$.

For modules use the reduced bar comodule
$T^c(s\overline{\mathcal K})\otimes M$, where
$\overline{\mathcal K}=\mathbb Q\xi$ and $s$ now lowers resolution
degree by one. An $A_\infty$ module structure on $M$ is a square-zero degree-one
differential on this comodule compatible with the algebra bar
differential. Here the reduced algebra bar differential is zero.

Fix the left-module sign so that the action term on
$(s\xi)\otimes a$ is $\xi a$. The operation $m_{j+1}$ has $j$
algebra inputs, one module input, and resolution degree $1-j$.
The adjective reduced excludes inputs equal to the algebra unit.

\begin{theorem}\label{df:thm:block}
The minimal strictly unital $A_\infty$ module of $M_{C_m}$ has exactly
one nonzero reduced operation:
\begin{equation}\label{df:eq:block-operation}
 m_{m+1}(\underbrace{\xi,\ldots,\xi}_{m},u)=v.
\end{equation}
All reduced operations on $v$ vanish. A bar comparison with $M_{C_m}$
has components
\begin{equation}\label{df:eq:block-comparison}
 I_j((s\xi)^{\otimes j},u)=(-1)^j e_{m-1-j}\quad(0\leq j<m),
 \qquad I_0(v)=(-1)^{m-1}\xi e_0,
\end{equation}
and all other components zero. In particular, the higher module
structure determines $m$.
\end{theorem}

\begin{proof}
The transferred component with $j$ algebra inputs is
$p\xi(-h\xi)^{j-1}i$. Each application of $h\xi$ lowers the basis
index by one. Projection first survives at $j=m$, when the sign
chosen for $v$ gives \eqref{df:eq:block-operation}. It vanishes for
larger $j$. Starting from $v$, the first exterior multiplication
already vanishes.

Since $s\xi$ has degree zero, the comparison equation on these words is
\[
 dI_j+\xi I_{j-1}
       =\sum_{k=1}^j I_{j-k}m_{k+1}.
\]
For $j<m$, the two terms on the left cancel. At $j=m$, their sum
is $I_0(v)$. Above $m$, both sides vanish. The minimal bar differential
squares to zero since its only nonzero output is $v$, on which every
reduced operation is zero. Strict units give the remaining identities.

The coefficient dual to the Taylor map in
\eqref{df:eq:block-operation} is $z^m$. The corresponding two-term
free presentation is $[R_0\xrightarrow{z^m}R_0]$, up to an overall
shift. Its only nonzero cohomology is $R_0/(z^m)$. Deleting the higher
map would change this module while leaving the two cohomology
dimensions of $M_{C_m}$ equal to one.
\end{proof}

For $m=1$ the operation is the ordinary binary action. For $m>1$ the
binary action on cohomology is zero, and the first nonzero reduced
operation has arity $m+1$. If $e_0$ has internal degree $k$, then
$u$ has internal degree $k+2m-2$ and $v$ has degree $k-2$.
Equation~\eqref{df:eq:block-operation} preserves that grading.

Apply the theorem to $C=H^*(\mathbf P(TX);\mathbb Q)$, with
$N=c_1(L)\cup-$. The Jordan lengths computed in
Chapter~\ref{ed:chap:euler} (Proposition~\ref{ed:prop:jordan} and
Theorem~\ref{ed:thm:square}) give the following multiplicities:
\begin{equation}\label{df:eq:module-table}
 \begin{array}{c|rrrr}
  X&m_3&m_5&m_7&m_9\\ \hline
  \text{Enriques surface}&10&1&0&0\\
  \text{K3 surface}&22&1&0&0\\
  \text{K3 Hilbert square}&23&276&23&1
 \end{array}.
\end{equation}
For the Hilbert square there are $323$ blocks, but their lengths sum
to $1296$. The two resolution cohomology dimensions each see only
$323$. The higher module operations distinguish all four lengths.
Although their displayed decomposition uses a Jordan basis,
the dg module \eqref{df:eq:parameter-module} is defined directly
by the tautological operator.

This construction also retains maps. If $T:C\to C'$ satisfies
$TN=N'T$, then $T\otimes1$ is a strict dg $\mathcal K$-module map.
Choose contractions and transfer them to the bar modules, writing
$I,P,H$ for the resulting maps with $IP-1=-dH-Hd$.
The transferred map is $F(T)=P'TI$.

Let $T_1:C_1\to C_2$ and $T_2:C_2\to C_3$ commute with the respective
parameter operators. Subscripts on $I,P,H$ label these modules. Then
\begin{equation}\label{df:eq:composition-homotopy}
 F(T_2)F(T_1)-F(T_2T_1)=dK+Kd,
 \qquad K=-P_3T_2H_2T_1I_1.
\end{equation}
This follows by substituting $I_2P_2-1=-dH_2-H_2d$.
The perturbation formulas are finite at each bar length, since each
action lowers that length. Thus changes of contraction retain the
module maps and their composition homotopies.

For comparison, a free resolution of a cyclic block gives
\begin{equation}\label{df:eq:block-hom}
 \begin{split}
 \operatorname{Hom}_{R_0}(C_m,C_n)
   &=z^{\max(n-m,0)}R_0/(z^n),\\
 \operatorname{Ext}_{R_0}^1(C_m,C_n)
   &=R_0/(z^{\min(m,n)}),\qquad
 \operatorname{Ext}_{R_0}^{j}(C_m,C_n)=0\quad(j>1).
 \end{split}
\end{equation}
Indeed, applying Hom to
$[R_0\xrightarrow{z^m}R_0]\to C_m$ gives the kernel and cokernel
of $z^m$ on $C_n$.

Finite-length $R_0$-modules form the extension closure of $\mathbb Q$.
Derived Hom from $\mathbb Q$ sends this generator to the free
$\mathcal K$-module and identifies their derived endomorphisms.
Exact triangles extend the identification to all finite-length modules.
The minimal modules therefore retain the groups in \eqref{df:eq:block-hom}.

\section{Split embeddings and Gysin maps}
\label{df:sec:gysin}

Return to cohomological grading. Let $E$ and $F$ be complex bundles
over $X$, of ranks $r\geq1$ and $s\geq0$. Choose Chern representatives
from direct-sum connections, so that
$e_z(E\oplus F)=e_z(E)e_z(F)$ holds as a cochain identity.
For the inclusion $j:\mathbf P(E)\hookrightarrow\mathbf P(E\oplus F)$,
the normal bundle is $L^*\otimes\pi^*F$. Its Euler class is $e_z(F)$
after identifying $z$ with $c_1(L)$.

\begin{theorem}\label{df:thm:gysin}
The pullback and the degree-$2s$ Gysin map admit the following
$A_X[z]$-linear representatives on Euler cochain models:
\begin{equation}\label{df:eq:gysin}
 \begin{aligned}
 j^*(a)&=a,&
 j^*(\eta_{E\oplus F}b)&=\eta_E e_z(F)b,\\
 j_!(a)&=e_z(F)a,&
 j_!(\eta_E b)&=\eta_{E\oplus F}b.
 \end{aligned}
\end{equation}
They satisfy the projection formula, the self-intersection formula
$j^*j_!=e_z(F)$, and composition for flags of split bundles.
\end{theorem}

\begin{proof}
Write odd generators on the left. For homogeneous $b$,
$D(\eta_E b)=e_z(E)b-\eta_E Db$. Substitution into
\eqref{df:eq:gysin} proves the chain-map identities, including
coefficients of odd degree or nonzero differential. The same formulas
give $j_!(xj^*y)=j_!(x)y$ and $j^*j_!(x)=e_z(F)x$.
For $E\subset E\oplus F\subset E\oplus F\oplus G$,
the even factors multiply as $e_z(F)e_z(G)$, while the odd generator
is sent to the final odd generator. This proves flag compatibility.

For pullback, retain the fixed connection on $L^*\otimes\pi^*F$.
On $\pi^*E\otimes L^*$, choose the connection $\nabla_0$ trivial
on the tautological section. Their direct sum preserves the
tautological section of the restricted bundle.
The top Chern transgression is then $\Theta_E e_\lambda(F)$.
This realizes the first line of \eqref{df:eq:gysin} geometrically.

For the Gysin map, first work over the universal base
$BU(r)\times BU(s)$. Its rational cohomology is the polynomial
algebra in the Chern classes of $E$ and $F$, all in even degrees.
Together with $z$, call this algebra $R$. The two projective-bundle
models resolve $R/(e_z(E))$ and $R/(e_z(E)e_z(F))$.
The projection $L\to\pi^*F$ cuts out $\mathbf P(E)$ transversely.
Its Euler class gives $j_!(1)=e_z(F)$.

A degree-$2s$ $R$-linear map from the free Euler resolution of
$R/(e_z(E))$ to $R/(e_z(E)e_z(F))$ is determined by its value on $1$.
Its value on $\eta_E$ would have odd degree $2(r+s)-1$, and the
target has no odd part. The relation is respected precisely because
$e_z(E)e_z(F)=0$ there. Thus the Gysin class has the representative
whose value on $1$ is $e_z(F)$. Lifting this representative to the
target Euler resolution gives exactly the second line of
\eqref{df:eq:gysin}.

The Thom construction commutes with pullback of the pair of bundles.
Derived base change from the universal, simply connected base therefore
gives the same Gysin map over $X$. This identifies its class in the
derived category of $A_X[z]$-modules.
\end{proof}

These formulas apply to a split normal bundle and its Thom class.
A Hall correspondence has a deformation complex involving cross
derived Hom terms. To use the same construction there, one must compare
that complex, its orientation, and its coefficient pullback and pushforward
with the Euler models. The split calculation specifies the required
compatibilities but does not supply this comparison.

\section{Maps of Euler cones and the vacuum}
\label{df:sec:cones}

The odd generators above also express a general condition on derived
quotient maps. Let $f:A\to B$ and $g:C\to D$ be degree-zero cochain
maps. Their cones have underlying spaces $B\oplus A[1]$ and
$D\oplus C[1]$, with differentials
\[
 d_f=\begin{pmatrix}d_B&f\\0&-d_A\end{pmatrix},\qquad
 d_g=\begin{pmatrix}d_D&g\\0&-d_C\end{pmatrix}.
\]
For multiplication by a homogeneous Euler class of degree $2r$,
the source shift $A[-2r]$ makes this convention applicable.

A degree-$k$ map preserving the target subcomplexes $B$ and $D$
has the form
\begin{equation}\label{df:eq:cone-map}
 T=\begin{pmatrix}b&h\\0&(-1)^k a\end{pmatrix},
 \qquad |a|=|b|=k,\quad |h|=k-1.
\end{equation}
These maps form a Hom complex with differential
$\partial T=d_gT-(-1)^kTd_f$. Direct multiplication gives
$\partial^2=0$ and the Leibniz rule for composition. In degree zero,
closedness says that $a$ and $b$ are chain maps and
\begin{equation}\label{df:eq:cone-homotopy}
 d_Dh+hd_A=bf-ga.
\end{equation}
The off-diagonal entry is therefore the specified homotopy making
the coefficient square commute. Composing two closed degree-zero maps
gives $h_{21}=b_2h_1+h_2a_1$.

There is a simple obstruction at the empty object. Use only resolution
grading for this example, with all coefficient modules in degree zero.
The empty Euler map is $f_0=1:R_0\to R_0$. For a connected
positive-dimensional smooth projective $Y$, the nonempty map is
$f_1=e_z(TY):H^*(Y)[[z]]\to H^*(Y)[[z]]$.
An ordinary creator with $b(1)=1$ would require
$1=e_z(TY)a(1)$ in \eqref{df:eq:cone-homotopy}.
This is impossible: reduction at $z=0$ and then to cohomological
degree zero sends the right side to zero and the left side to one.

The Euler-transported creator with $b(1)=e_z(TY)$ instead has the
strict lift $a(1)=1$, $h=0$. Its source cone is nevertheless
$\operatorname{Cone}(1)$, which is acyclic, so the induced derived
map is zero. This does not alter the image-module action constructed
in Section~\ref{ed:sec:transport} of Chapter~\ref{ed:chap:euler}.
It shows why that action and an
action on the Euler quotient require different vacuum data.

To identify either higher structure with chiral collision operations,
one must construct a map from the relevant supported operation complex.
It must preserve the intrinsic sphere-fiber products or the parameter
module operations on their respective carriers, together with the
vacuum comparisons.
